\subsection{讨论}\label{sec:5.9}

\noindent\textbf{\autoref{sec:5.1}.} 关于度量导数与连续性方程的详细论述，见\cite{ref18,ref281,ref317}. 特别地，\autoref{thm:5.3}的严格版本见\cite[第 8 章]{ref18}或\cite[第 5 章]{ref281}.

\noindent\textbf{\autoref{sec:5.2}.} 有许多介绍黎曼几何的优秀教材，例如\cite{ref113}.

\noindent\textbf{\autoref{sec:5.3}.} Benamou--Brenier 公式常称为最优传输的“动力学”表述，与描述“静态”图景的\autoref{sec:1}相对. 对偶问题也有动力学版本，其中对偶势按照 Hamilton--Jacobi 方程演化，见\cite[第 8.1 节]{ref315}. \autoref{sec:4}引入的熵正则最优传输的动力学版本，与著名的 Schrödinger 桥问题密切相关\cite{ref70,ref206}.

\noindent\textbf{\autoref{sec:5.4}.} 本节的形式计算法则最初由 Otto\cite{ref249}提出，不过更早的工作\cite{ref191}已经预示了其中一些思想.

测度 $\mu$ 处的切空间连同其度量，可以视为 $\mu$ 处几何结构的线性化. 由此得到“线性化最优传输”，成为许多应用的基础\cite{ref33,ref35,ref63,ref180,ref181,ref255,ref284,ref322}.

除了梯度流，也有人提出将其他优化算法适配到 Wasserstein 空间，例如\cite{ref73,ref88,ref301,ref320,ref321}.

\noindent\textbf{\autoref{sec:5.5}.} Bures--Wasserstein 几何以 Donald Bures\cite{ref62}命名；他在量子信息论研究中，在半正定锥上引入了这一度量. 关于 BW 几何的进一步研究，见\cite{ref34,ref162,ref230,ref311}. \autoref{exer:5.12}描述的与 Burer--Monteiro 分解的联系，已在低秩矩阵恢复背景下研究\cite{ref218,ref224}.

\noindent\textbf{\autoref{sec:5.6}.} 正文所述的本节几何最初在\cite{ref69,ref107}中考虑，而梯度流方程由\cite{ref192}得到.

\noindent\textbf{\autoref{sec:5.7}.} Fisher--Rao 几何在信息几何中已有深入研究\cite{ref16,ref26}.

\noindent\textbf{\autoref{sec:5.8}.} 一些来源，例如\cite{ref209}，用权重的平方根 $r:=\sqrt w$ 表示 WFR 梯度流. 这一约定的动机是：FR 距离对应于权重平方根之间的欧氏距离，因此可以将 WFR 距离解释为赋予“锥”度量的 $(r,x)$ 对空间上的耦合代价（参见\autoref{sec:7.3.3}）. 因为这里不介绍这一视角，所以采用更直接的、以 $w$ 为变量的参数化.

高斯粒子的使用最初由\cite{ref192}倡导.
